% ---------- ¶1  puzzle, then the theoretical benchmark ----------
Mandatory disclosure is the standard regulatory response to asymmetric information about
product quality, yet theory suggests markets may not need it. When quality is verifiable, the
unraveling result of \citet{grossman_informational_1981} and \citet{milgrom_good_1981} holds that sellers disclose
voluntarily, since buyers rationally read silence as the worst possible news. All but the very
worst sellers therefore disclose, and disclosure becomes effectively universal without a
mandate. Whether this result describes real markets remains unsettled: it rests on strong
assumptions and is difficult to test, because in most settings where quality is verifiable
disclosure has long been compulsory.
